\documentclass[11pt]{article}
\usepackage{amsmath,amsthm,amssymb,graphicx,xcolor}
\usepackage[T1]{fontenc}
\usepackage{lmodern}
\usepackage[colorlinks=true]{hyperref}
\usepackage{tikz}
\usepackage{ifthen}
\newcommand{\expref}[2]{\texorpdfstring{\hyperref[#2]{#1~\ref{#2}}}{#1~\ref{#2}}} 
\newtheorem{theorem}{Theorem}[section]
\newtheorem{lemma}[theorem]{Lemma}
\newtheorem{corollary}[theorem]{Corollary}
\newtheorem{claim}[theorem]{Claim}
\newtheorem{proposition}[theorem]{Proposition}
\newtheorem{fact}[theorem]{Fact}
\newtheorem{conjecture}[theorem]{Conjecture}
\newtheorem{definition}[theorem]{Definition}
\newtheorem{example}[theorem]{Example}
\newtheorem{remark}[theorem]{Remark}
\newtheorem{acknowledgement}[theorem]{Acknowledgement}
\newtheorem{axiom}[theorem]{Axiom}
\newtheorem{case}[theorem]{Case}
\newtheorem{conclusion}[theorem]{Conclusion}
\newtheorem{condition}[theorem]{Condition}
\newtheorem{criterion}[theorem]{Criterion}
\newtheorem{exercise}[theorem]{Exercise}
\newtheorem{notation}[theorem]{Notation}
\newtheorem{problem}[theorem]{Problem}
\newtheorem{question}[theorem]{Question}
\newtheorem{solution}[theorem]{Solution}
\newtheorem{summary}[theorem]{Summary}
\newcommand{\R}{{\mathbb R}}
\newcommand{\eee}{\mathrm e}
\newcommand{\ie}{i.\,e.} 
\newcommand{\eg}{e.\,g.}
\newcommand{\Eg}{E.\,g.}
\usepackage{packages}
\usepackage{aumacros}
\usepackage{eprint}

\usepackage[margin=25mm]{geometry}
\emergencystretch=3em
\providecommand{\Z}{\mathbb{Z}}
\providecommand{\N}{\mathbb{N}}
\providecommand{\Q}{\mathbb{Q}}
\providecommand{\C}{\mathbb{C}}
\providecommand{\Ie}{I.\,e.}
\usepackage{fontspec}
\newfontfamily\collectioncyrillicfont{DejaVu Serif}
\renewcommand{\foreignlanguage}[2]{{\collectioncyrillicfont #2}}
\title{Local-to-global spectral descent for weighted simplicial complexes}
\author{Prahladh Harsha, Ramprasad Saptharishi}
\date{doi:10.4086/toc.2024.v020a005; CC BY 3.0}
\begin{document}
\maketitle
\noindent\textit{Editorial note.} Complete original mathematical body and bibliography retained. General descent theorem counted once, with full original Oppenheim attribution and Dikstein-exposition credit; concrete coset construction and one-level supporting theorem remain context without additional counting in this task.\par\medskip
\noindent\textit{Editorial note.} Unused nth ordinal-formatting import omitted in render copy; original package list preserved.\par\medskip
\noindent\textit{Editorial note.} Standard number-set macros supplied from the original ToC convention, with wider article text area for original long formulas. No mathematical edits.\par\medskip
\noindent\textit{Editorial note.} Supplied original ToC prose abbreviation Ie in article wrapper.\par\medskip
\noindent\textit{Editorial note.} Original Cyrillic bibliography text retained with existing DejaVu Serif through a local foreignlanguage presentation wrapper; no transliteration, global language installation or mathematical edits.\par\medskip
\noindent\textit{Allowed standard prerequisites (editorial).}\begin{itemize}
\item Self-adjoint spectral operators, weighted simplicial links and conditional expectations. All local weighted identities, the one-level argument and dimensional iteration are retained. Original attribution to Oppenheim and exposition from Dikstein remain intact.
\end{itemize}\medskip
\section{Introduction}

In the last few years, there has been a surge of activity related to high-dimensional expanders (HDXs).
%
High-dimensional  %
expanders are
%
%
high-dimensional generalizations   %
of classical graph expanders.
Depending on which definition of graph expansion is generalized, 
there are several different (and unfortunately, 
%
often  %
%
%
%
inequivalent) definitions of HDXs.
For the %
purposes  %
of this note, we will restrict ourselves to the spectral definition 
of HDXs (see \expref{Definition}{def:hdx}).
%
\emph{Ramanujan graphs} are families of expander graphs which have optimal spectral expansion (see \expref{Definition}{def:spectral-expander}). Lubotzky,
Phillips and Sarnak \cite{LubotzkyPS1988} and Margulis \cite{margulis88} independently gave explicit 
constructions of %
Ramanujan graphs. \emph{Ramanujan complexes} are high-dimensional generalizations of Ramanujan graphs. 
Lubotzky, Samuels and Vishne~\cite{LubotzkySV2005-exphdx,LubotzkySV2005-hdx} and 
Sarveniazi \cite{Sarveniazi2007}
gave explicit constructions of %
 Ramanujan complexes, which are the high-dimensional analgoues
of the construction %
of Lubotzky, %
Phillips  %
and Sarnak~\cite{LubotzkyPS1988}.   %
%
%
%
%
%
%
%
%
%
These constructions were  %
the first construction of constant-degree %
spectral HDXs.
The Ramanujan %
graphs constructed by Lubotzky,  %
Phillips and Sarnak \cite{LubotzkyPS1988} and Margulis \cite{margulis88} have the 
nice property that they are simple to describe, 
however %
the proof of the optimality of their %
%
expansion is involved.
%
%
%
%
%
The Ramanujan %
complexes constructed by  %
Lubotzky, Samuels and Vishne~\cite{LubotzkySV2005-exphdx,LubotzkySV2005-hdx} and 
Sarveniazi \cite{Sarveniazi2007}, on the other hand, are non-trivial to describe
and it is difficult to prove their high-dimensional expansion property. %
%
Subsequently Kaufman and Oppenheim~\cite[See STOC version]{KaufmanO2018} gave an
%
%
elegant 
%
elementary construction of bounded-degree %
spectral HDXs using elementary matrices.
%
%
%
%
%
%
While their HDXs are not Ramanujan, their construction gives rise to
new families of expander graphs whose spectral gap is close to the
optimal Ramanujan bound.   %
%
%

Despite %
the Kaufman--Oppenheim   %
construction being elementary and simple, the proof of expansion
%
is  %
not elementary. 
%
%
%
%
%
%
%
%
%
%
%
%
%
%
%
%
%
%
The purpose of this exposition is to give an elementary proof 
%
of the expansion of the Kaufman--Oppenheim HDX construction. 
In particular, we obtain the same eigenvalue bound as their proof.  
%
%
%
%
%

%
The 1-skeleton (see \expref{Def.}{def:simplicial-complex}) %
of a HDX %
(even of a one-sided spectral HDX) is a 
two-sided spectral expander (see \expref{Def.}{def:spectral-expander}). %
Thus, this construction has the added advantage that it yields an elementary 
construction (accompanied by a simple proof) of a standard two-sided spectral
expander (though not an optimal one).  %
%
%
%

%

\section{Preliminaries}

We begin by recalling what a simplicial complex is. 
\begin{definition}[Simplicial complex]   \label{def:simplicial-complex}
  A \emph{simplicial complex} $X$ over a finite set $U$ is a collection of subsets of $U$ with the property that if $S \in X$ then any $T \subseteq S$ is also in $X$.
  \begin{itemize}
  \item For all $i\geq -1$,
    define $X(i) := \setdef{S \in X}{|S| = i+1}$. Thus, if $X$ is
    non-empty, then $X(-1) = \set{\emptyset}$. 
  \item The elements of $X$ are called \emph{simplices} or
    \emph{faces}. The elements of $X(0)$, $X(1)$ and $X(2)$ are
    usually referred to as \emph{vertices}, \emph{edges} and
    \emph{triangles}, respectively.  %
%
%
%
%
%
%
    \item For any $1 \leq k \leq d$, the \emph{$k$-skeleton} of the complex
    $X$ is the subcomplex $X(-1) \cup X(0) \cup X(1) \cup \dots \cup X(k)$.
      We identify the \emph{1-skeleton} with the graph defined by $X(0)$ and
      $X(1)$.   %
  \item The \emph{dimension} of the simplicial complex $X$ is defined as the largest $d$ such that $X(d)$ (which consists of faces of size $d+1$) is non-empty.
  \item The simplicial complex is said to be \emph{pure} if every face
    is contained in some face in $X(d)$, where $d = \dim(X)$.
    
  \item For a face $S \in X$, the \emph{link} of $S$, denoted by
    $X_S$, is the simplicial complex defined as 
    \[
      X_S := \setdef{T \setminus S}{T\in X\;,\; S \subseteq
        T}. \qedhere\
    \]
  \end{itemize}
\end{definition}

Thus, a graph $G=(V,E)$ is just a simplicial complex $G$ of dimension one
with $G(0) = V$ and $G(1)=E$. We will deal with \emph{weighted} pure
simplicial complexes where the weight function satisfies a certain
\emph{balance} condition.

\begin{definition}[weighted pure simplicial complexes]
Given a $d$-dimensional pure simplicial complex $X$ and an associated
weight function $w\colon X \to \R_{\geq 0}$, we say the weight
function is \emph{balanced} if the following two conditions are
satisfied.
\begin{align}\label{eq:weight}
  \sum_{S \in X(d)} w(S) & = 1 \; ; & 
  w(S) &  = \frac{1}{i+2}\sum_{T \in X(i+1), T \supset S}
              w(T), \ \text{ for all } i<d \text{ and } S \in X(i).
\end{align}
%
%
%
%
A \emph{weighted simplicial complex} $(X,w)$ is a pure simplical
      complex accompanied with a balanced weight function $w$. If no
      weight function is specified, then we work with the balanced weight
      function $w$ induced by the uniform distribution on the set
      $X(d)$ of maximal faces.

      For a face $S \in X$, the balanced weight function $w_S$ associated
      with the link $X_S$ is the restricted weight function, suitably
      normalized, more precisely $w_S :=
    \sfrac{w|_{X_S}}{w(S)}$. 
\end{definition}

Condition~\eqref{eq:weight} states that the weight function can be interpreted as a family of joint distributions $(w|_{X(-1)}, \ldots, w|_{X(d)})$ where $w|_{X(i)}$ is a probability distribution on $X(i)$.
The distribution $w|_{X(d)}$ is specified by the first condition in \eqref{eq:weight} while the second condition implies that the weight distribution $w|_{X(i)}$ is the distribution on $X(i)$ obtained by picking a random $\tau \in X(d)$ according to $w|_{X(d)}$ and then removing $(d-i)$ elements uniformly at random.


We now recall the classical definition
of what it means for a graph to be a spectral expander.
We will be exclusively discussing only \textbf{undirected graphs.}
%
%
%

%
%

\begin{definition}[spectral expander]  \label{def:spectral-expander}
Given %
a %
weighted graph $G=(V,E, w)$ on $n$ vertices, let $A_G$ be its
normalized adjacency matrix given as follows:
\[ A_G(u,v) := \begin{cases}
        \frac{w(u,v)}{w(u)} & \text{if } \{u,v\} \in E,\\
      0 & \text{otherwise.}
    \end{cases}
    \]
Let $1 = \lambda_1 \geq \lambda_2 \geq
\cdots \geq \lambda_n \geq -1$ be the $n$ eigenvalues%
%
%
%
%
%
%
%
%
\footnote{By the balance condition, $w$ satisfies
  $w(v) = \sum_{\{u,v\} \in E} w(u,v)$. The
  matrix $A_G$ is self-adjoint with respect to the
  inner product $\inangle{f,g}_w := \E_{v \sim w}[f(v) g(v)]$ since
  $\inangle{f,Ag}_w = \inangle{Af,g}_w = \E_{\{u,v\} \sim w}[f(u)
  g(v)]$. Hence, $A_G$ has $n$ real eigenvalues which can be
  obtained using the %
  Courant--Fischer Theorem (\expref{Theorem}{thm:cf1}).}
  %
 of $A_G$.  We denote the second largest
eigenvalue of $G$ %
by  %
$\lambda(G)$. 

$G$ is said to be a
\emph{$\lambda$-spectral expander} if $\max\{\lambda_2,|\lambda_n|\}
\leq \lambda$. This is sometimes also referred to as a
$\lambda$-two-sided spectral expander.

$G$ is said to be a
\emph{$\lambda$-one-sided spectral expander}
if $\lambda_2 \leq \lambda$.
\end{definition}

This spectral definition of expanders is generalized to
higher-dimensional  %
simplicial complexes as follows.

\begin{definition}[$\lambda$-spectral HDX]\label{def:hdx}
  A weighted simplicial complex $(X,w)$ of dimension $d \geq 1$ is said to be a
  \emph{$\lambda$-spectral HDX} (or a
\emph{$\lambda$-two-sided spectral HDX})\footnote{These 
     are sometimes also referred
    to as $\lambda$-link HDXs or $\lambda$-local-expanders to
    distinguish from an alternative global definition of
    high-dimensional expansion.} if for every $-1\leq i \leq d-2$ and
  $s \in X(i)$, the weighted $1$-skeleton of the link $(X_s, w_s)$ is a
  $\lambda$-spectral expander.

A weighted simplicial complex $(X,w)$ of dimension $d \geq 1$ is said to be a
 \emph{$\lambda$-one-sided spectral HDX}
if for every $-1\leq i \leq d-2$ and
  $s \in X(i)$, the weighted $1$-skeleton of the link $(X_s,w_s)$ is a
  $\lambda$-one-sided spectral expander.
\end{definition}

Using Garland's technique~\cite{Garland1973}, Oppenheim~\cite{Oppenheim2018} 
showed that if the 1-skeletons of all the links are
connected, then a spectral gap at dimension $(d-2)$ descends to all
lower levels. 

%
\begin{dtheorem}[Oppenheim, 2018]\label{thm:trickle-down} %
%
%
%
Let $(X,w)$ be  %
a $d$-dimensional weighted simplicial complex with the following properties.
  \begin{itemize}
  \item For all $s \in X(d-2)$, the link $(X_s,w_s)$ is a
    $\lambda$-one-sided spectral expander for some $\lambda < \frac{1}{d-1}$. 
  \item The $1$-skeleton of every link is connected.
  \end{itemize}
  Then, $(X,w)$ is a $\pfrac{\lambda}{1 - (d-1)\lambda}$-one-sided spectral HDX.
\end{dtheorem}
Thus to prove that a given simplicial complex is a spectral HDX, it suffices to show that the 1-skeleton of %
every link is %
connected and %
to show   %
a spectral gap at the top level.
For the sake of completeness, we give a proof of the 
%
Descent Theorem  %
in
%
\expref{Appendix}{sec:descent}  %
which includes a descent theorem for the least eigenvalue as well.

\section{Coset complexes}

The HDX construction of Kaufmann and Oppenheim is a particular instantiation of a certain type of simplicial complex called a \emph{coset complex} based on a group and its subgroups. In this section, we give an exposition of these objects. 
For a basic primer on group theory,
see \expref{Section}{sec:group}.  %
%

\begin{definition}[coset complex] Let $G$ be a group and let
  $K_1,\ldots, K_d$ be $d$ subgroups of $G$. The
coset complex $\X(G,\set{K_1,\ldots, K_d})$ is a $(d-1)$-dimensional
simplicial complex defined as follows:
\begin{itemize}
\item The set of \emph{vertices}, $\X(0)$, 
consists of
the %
%
%
left cosets of $K_1,\ldots, K_d$ and we shall say
the %
left  %
cosets of $K_i$ are of \emph{type} $i$.
\item The set of \emph{maximal faces}, $\X({d-1})$,
consists of the $d$-sets of
cosets of different types with a non-empty intersection. That is,
  \[
    \set{g_1 K_1,\ldots, g_d K_d}\in \X(d-1) \Longleftrightarrow g_1 K_1 \cap \cdots \cap g_d K_d \neq \emptyset.
  \]
  An equivalent way of stating this is that $\set{g_1 K_1,\ldots, g_d K_d}\in \X(d-1)$ if and only if there is some $g \in G$ such that $g_i K_i = g K_i$ for all $i$, since
  \[
    g_i K_i = g K_i \;\Longleftrightarrow\; K_i = g_i^{-1} g K_i  \;\Longleftrightarrow\; g_{i}^{-1} g \in K_i \;\Longleftrightarrow\; g \in g_{i} K_i.
  \]
\item The 
lower-dimensional %
faces are obtained by \emph{down-closing}
  the maximal faces. Hence, for $0\leq r \leq d$, $\set{g_{i_1} K_{i_1}, \ldots, g_{i_r} K_{i_r}} \in \X(r-1)$ if and only if $i_j \neq i_k$ for all $j \neq k$ and
  \[
    g_{i_1} K_{i_1} \cap \cdots \cap g_{i_r} K_{i_r} \neq \emptyset.
  \]
  We shall call the set $\set{i_1,\ldots, i_r}$ the \emph{type} of this face.
\item The dimension of this complex is $d-1$.
\item The weight function we will use is the one induced by the
  uniform distribution on the set $\X(d-1)$ of maximal faces. 
  \qedhere
\end{itemize}
\end{definition}

A simplicial complex constructed this way is \emph{partite} in the sense that each maximal face consists of vertices of distinct types. 

It follows from the definition, that $\X(i)$ is precisely the set of
cosets of the form $g K_S$ where $K_S = \Intersection_{j\in S} K_j$
for sets $S \subseteq [d]$ of size exactly $i+1$. In particular, $\X(d-1)$, the
set of maximal faces, is in 1-1 correspondence with the group
$G$ if $\Intersection_{j\in [d]} K_j = \{\mathrm{id}\}$ where
``$\mathrm{id}$'' is the identity element of the group $G$. 



\subsubsection*{Connectivity}  %

\begin{observation}\label{obs:nonempty-coset-intersection}
$g_1 K_1 \cap g_2 K_2 \neq \emptyset$ if and only if $g_1^{-1} g_2 \in K_1 K_2$. 
\end{observation}
\begin{proof}
  ($\Rightarrow$) Say $x = g_1 k_1 = g_2 k_2$ for $k_1 \in K_1$ and $k_2 \in K_2$. Then $g_1^{-1} g_2 = g_1^{-1} x \cdot x^{-1} g_2 = k_1 k_2^{-1} \in K_1 K_2$.

  ($\Leftarrow$) If $g_1^{-1} g_2 = k_1 k_2$ for
some  %
$k_1 \in K_1$ and $k_2 \in K_2$, then $g_1 k_1 = g_2 k_2^{-1} \in g_1 K_1 \cap g_2 K_2$. 
\end{proof}

\begin{lemma}[Criterion for connected $1$-skeletons]\label{lem:conn-criterion}
The $1$-skeleton %
%
%
%
%
of $\X(G,\set{K_1,\ldots, K_d})$ is connected if and only if 
$G = \inangle{K_1,\ldots, K_d}$. 
\end{lemma}
\begin{proof}
  ($\Leftarrow$) Since there is always an edge between $gK_i$ and $gK_j$ for $i\neq j$, it suffices to show that $K_1$ is connected to $gK_1$ for an arbitrary $g \in G$.
Suppose, for an arbitrary element $g \in G$, we have $g = g_1 \ldots
g_r$ where $g_j \in K_{i_j}$ and $i_j \neq i_{j+1}$ for each $j$. We might, without loss of generality, assume that
(a) $g_1 \in K_1$ (otherwise set $g = 1\cdot g_1 \cdots g_r$)  and (b) if $r \geq 2$,
then $i_r \neq 1$ (since otherwise we might then have worked with $g' =
g_1g_2 \ldots g_{r-1}$ as $gK_1= g'g_r K_1 = g'K_1$).

Then, we get the following path connecting $K_1 $ and $g K_{i_r}$
  \[
    K_1 =  g_1 K_{i_1} \rightarrow (g_1 g_2) K_{i_2} \rightarrow (g_1 g_2 g_3) K_{i_3} \rightarrow \ldots \rightarrow (g_1\cdots g_r) K_{i_r} = g K_{i_r}.
  \]
  Note that, due to \expref{Observation}{obs:nonempty-coset-intersection}, each
  successive pair of cosets are connected by an edge in the simplicial
  complex. Now, since $gK_{i_r}$ is adjacent to $gK_1$ (as $i_r \neq
  1$), we have that $K_1$ is connected to $gK_1$. \\

  \noindent
  ($\Rightarrow$) For an arbitrary $g \in G$, since the $1$-skeleton is connected we have a path
  \[
    K_1 = g_0 K_{i_0} \rightarrow  g_1 K_{i_1} \rightarrow \cdots \rightarrow g_r K_{i_r} = g K_1.
  \]
  By \expref{Observation}{obs:nonempty-coset-intersection}, for every $j = 0,\ldots, r-1$, we have $g_j^{-1} g_{j+1} \in K_{i_j} K_{i_{j+1}} \in \inangle{K_1,\ldots, K_d}$.
  Therefore,
  \[
    g = (g_0^{-1} g_1) \cdot (g_1^{-1} g_2) \cdots (g_{r-1}^{-1} g_r) \in \inangle{K_1,\ldots, K_d}. \qedhere
  \]  
\end{proof}


\subsubsection*{Structure of links of the coset complex}  %


For any set $S \subseteq [d]$, define the group $K_S :=
\Intersection_{i\in S} K_i$; let $K_\emptyset := \inangle{K_1,\ldots,
  K_d}$. The following lemma shows that the links of a coset complex
are themselves coset complexes. 

\begin{lemma}\label{lem:link-structure}
  For any $v\in \X(k)$ of type $S \subseteq [d]$, the link $X_v$ is isomorphic to the simplicial complex defined by $\X(K_S, \setdef{K_{S} \cap K_i}{i\notin S})$. 
\end{lemma}
\begin{proof}
  It suffices to prove this lemma for $v \in \X(0)$ as links of higher levels can be obtained by inductive applications of this case.

Observe that if $g$ is any element of  $G$, then
  $(g_{i_1} K_{i_1}, \ldots, g_{i_r} K_{i_r}) \in \X(r-1)$ if and only
  if $(g g_{i_1} K_{i_1}, \ldots, g g_{i_r} K_{i_r}) \in \X(r-1)$.
  Therefore, the link of a coset $g K_i$ is isomorphic to the link of
  the coset $K_i$. Thus, it suffices to prove the lemma for links of
  the type $X_{K_i}$ for some $i \in [d]$. 

  Let $v$ be the coset $K_1$, without loss of generality.
  The \emph{vertices} of the link, $X_v(0)$, are cosets of $K_2,\ldots, K_d$ that have a non-empty intersection with $K_1$.
  Note that any non-empty intersection $g_j K_j \cap K_1$ of a coset with $K_1$ is itself a coset $\tilde{g}_j (K_j \cap K_1)$ of the intersection subgroup $K_j \cap K_1$ in $K_1$. Indeed, suppose that $g_j h_j \in K_1$ for some $h_j \in K_j$. Then, $g_j h_j K_j = g_j K_j$ and $g_j h_j K_1 = K_1$ and hence
  \[
    g_j K_j \cap K_1 = g_j h_j K_j \cap g_j h_j K_1 = g_j h_j (K_j \cap K_1).
  \]
%
%
%
%
Therefore,\ the\ vertices\ of\ the\ link\ $X_v(0)$\ are\ in\ bijective\
correspondence\ with\
the\ %
%
cosets\ of\
the\ subgroups\ %
%
$\setdef{K_j \cap K_1}{j \in \set{2,\ldots, d}}$.

  The maximal faces in $X$ that contain the coset $K_1$ are precisely
the %
$d$-sets %
%
$\set{K_1, g_2 K_2, \ldots, g_d K_d}$
of cosets %
%
%
with a non-empty intersection and hence
  \[
    \emptyset \neq K_1 \cap g_2 K_2 \cap  \cdots \cap g_d K_d = (g_2 K_2 \cap K_1) \cap \cdots \cap (g_d K_d \cap K_1) = \tilde{g}_2 (K_2 \cap K_1) \cap \cdots \cap \tilde{g}_d (K_d \cap K_1),
  \]
  which are  precisely the maximal faces of the coset complex
  $\X\inparen{K_1,\setdef{K_j \cap K_1}{j \in \set{2,\ldots,
        d}}}$.  This establishes the isomorphism between $X_v$ and $\X\inparen{K_1,\setdef{K_j \cap K_1}{j \in \set{2,\ldots,
        d}}}$.
\end{proof}

\section{A concrete instantiation}

The simplicial complex of Kaufman and Oppenheim~\cite[See STOC version]{KaufmanO2018} is
a specific instantiation of the above \emph{coset complex}
construction. This section is devoted to an exposition of this
instantiation of Kaufman and Oppenheim.  We will need some notation to
describe their group.


\subsection*{Notation}

\begin{itemize}
\item %
  %
  %
  Let $p$ be a prime power and consider the ring $\F_p[t]$ of polynomials
  over the finite field $\F_p$.  %
  Let $R$ denote the
  quotient %
  ring %
  %
  %
  $\F_p[t]/\inangle{t^s}$  %
  where $\inangle{t^s}$ denotes the ideal generated by $t^s$.  %
  This is a ring whose elements can be identified with polynomials in
  $\F_p[t]$ of degree less than $s$ (where addition and multiplication
  are performed modulo $t^s$), with 
  %
  %
%
%
%
%
%
$t$ being a formal variable and $s$ a positive integer%
%
 \footnote{In the subsequent journal version \cite{KaufmanO2018}, the authors generalise
 the ring $R$ to any unital commutative ring. For the %
purposes  %
of this exposition, we focus exclusively on
finite  %
rings of the form described above.}.
%
%
%
%
%
By increasing the value of $s$, the construction will provide a family of complexes
on
a %
growing number of vertices. 
%


\item For any $d\geq 3$, and $1\leq i,j\leq d$ with $i \neq j$ and an
 element $r \in R$,
  we define  $e_{i,j}(r)$ to be the $d\times d$ elementary matrix with $1$'s 
on the diagonal and $r$ on the $(i,j)$-th entry.

For the sake of notational convenience, we shall
%
%
extend this notation 
  and %
write  %
%
%
$e_{k,\ell}(r):=e_{i,j}(r)$ for all $k,\ell\in\Z$ such that 
$k\equiv i\pmod{d}$ and $\ell\equiv j\pmod{d}$ $(1\le i,j\le d)$. %
%
For example, $e_{d,d+1}(r)$ refers to $e_{d,1}(r)$.

%
We will extend further and use $\set{i,i+1, \ldots, j-1}_{\bmod d}$ to denote
the set ${i,i+1, \cdots j-1}$ when $i < j$, and the set 
$\set{i, i+1, \ldots d, 1, 2, \ldots, j-1}$ when $j \leq i$. 
\end{itemize}

\noindent
We are now ready to describe the groups in the construction.
\begin{align*}
  \text{For $i\in \set{1,\ldots, d}$,}\quad K_i & = \inangle{e_{j,j+1}(at + b)\;:\; a,b \in \F_p\,,\,j \in [d] \setminus \set{i}}.\\
  G & = \inangle{K_1,\ldots, K_d}
\end{align*}
Each $K_i$ is generated by elementary matrices that have $1$'s on the
diagonal and an arbitrary linear polynomial in one entry of the
generalised diagonal %
$\set{(i,j)\;:\; i+1\equiv j \pmod d}$.  %

It so happens that the group $G$ generated by the subgroups
$K_1, \ldots, K_d$ is $\mathrm{SL}_d(R)$, the group of $d\times d$
matrices with entries in $R$ whose determinant is $1$ (in $R$). This
%
is a non-trivial fact %
(see  %
\cite[Theorem 4.3.9]{HM89}).   %
All we will need is the simpler fact that $|G|$
grows exponentially with $s$ (for fixed $p$ and $d$) while the size of
the groups $K_i$ are functions of $p$ and $d$ (and independent of
$s$). This will follow from the sequence of observations and lemmas
developed in the following section.\\

\noindent Given the above definition, there are two ``different''
subgroups we can define. 
\begin{align*}
K_S & = \Intersection_{i\in S} K_i,\\
\widetilde{K_S} &:= \inangle{e_{i,i+1}(at+b)\;:\; a,b\in \F_p\,,\, i\notin S}.
\end{align*}
That is, $K_S$ is the intersection of the groups $\setdef{K_i}{i\in S}$, and
$\widetilde{K_S}$ is the group generated by the intersection of the
sets of %
generators of the
%
$K_i$.  %
Thus, clearly, $\widetilde{K_S} \subseteq K_S$. The following lemma shows that
in fact the two groups are identical. 


\begin{lemma}[Intersections of $K_i$'s]\label{lem:intersection-property-KS}
  For any $S \subseteq [d]$,
  \[
   \widetilde{K_{S}} = \inangle{e_{i,i+1}(at +b)\;:\; a,b\in \F_p\,,\, i\notin S}
    = \Intersection_{i\in S} K_i = K_S
  \]
  In other words, the group generated by the intersection of generators equals the group intersection. 
\end{lemma}

We will prove this lemma in the following section by giving an
explicit description of the groups that makes the above lemma
evident. An immediate consequence of this lemma is that $K_{[d]} =
\{\mathrm{id}\}$ and hence $\X(d-1)$, the set of maximal faces, is in
1-1 correspondence with the group~$G$.   %

\subsection{Explicit description of the groups}

The following is an easy consequence of the definition of
$e_{i,j}(r)$. Note that $e_{i,j}(r)$ is defined only if $i \neq j$. 

\begin{observation}\label{obs:comm-generators}

  \begin{enumerate}[(a)]
    \item\label{item:comm-gen-sum} \emph{Sum:} $e_{i,j}(r_1) \cdot e_{i,j}(r_2) =
      e_{i,j}(r_1 + r_2)$.

      As a corollary, $e_{i,j}(r)^{-1} = e_{i,j}(-r)$. 
\item\label{item:comm-gen-product} \emph{Product:} If $i \neq \ell$, the commutator\footnote{The commutator of two elements $g,h$,
 denoted by $[g,h]$, %
 is defined as $g^{-1} h^{-1} g h$. (\expref{Definition}{defn:commutator})
 %
 %
}
       $[e_{i,j}(r_1),e_{k,\ell}(r_2)]$ behaves as follows.
  \[
    [e_{i,j}(r_1),e_{k,\ell}(r_2)] = 
    \begin{cases}
      e_{i,\ell}(r_1r_2) & \text{if $j = k$},\\
      \mathrm{id} & \text{if $j \neq k$}.
    \end{cases}\qedhere
  \]
  \end{enumerate}
\end{observation}
\begin{proof}
  Let $\mu_{i,j}$ denote the matrix that has a $1$ at the $(i,j)$-th entry, and $0$ everywhere. Then, \eqref{item:comm-gen-sum} follows as
  \begin{align*}
    e_{i,j}(r_1) e_{i,j}(r_2) &= (I + \mu_{i,j} r_1) \cdot (I + \mu_{i,j}(r_2))\\
    & = I + \mu_{i,j} \cdot (r_1 + r_2)\qquad\text{(since $\mu_{i,j}^2 = 0$ when $i\neq j$)}.
  \end{align*}

  As for \eqref{item:comm-gen-product}, we follow along a similar calculation. Note that 
  \[
  \mu_{i,j} \cdot \mu_{k,\ell} = \begin{cases} 0 & \text{if $j\neq k$}\\ \mu_{i,\ell} & \text{if $j = k$}.\end{cases}
  \]
  Therefore, 
  \begin{align*}
    [e_{i,j}(r_1), e_{k,\ell}(r_2)] & = (I - \mu_{i,j} r_1)\cdot (I - \mu_{k,\ell}r_2) \cdot (I + \mu_{i,j} r_1)\cdot (I + \mu_{k,\ell}r_2).
  \end{align*}
  When $j = k$ (along with the assumption that $i\neq j$, $k\neq \ell$), this simplifies to
  \begin{align*}
    [e_{i,j}(r_1), e_{k,\ell}(r_2)] &= (I - \mu_{i,j}r_1 - \mu_{k,\ell}r_2 + \mu_{i,\ell} \cdot r_1r_2) \cdot (I + \mu_{i,j}r_1 + \mu_{k,\ell}r_2 + \mu_{i,\ell} \cdot r_1r_2)\\
    & = I + \mu_{i,j}(r_1 - r_1) + \mu_{k,\ell}(r_2 - r_2) + \mu_{i,\ell}(r_1r_2 + r_1r_2 - r_1r_2)\\
    & = I + r_1r_2 \mu_{i,\ell}.
  \end{align*}
  If $j \neq k$, then we get
  \begin{align*}
  [e_{i,j}(r_1), e_{k,\ell}(r_2)] & = (I - \mu_{i,j} r_1)\cdot (I - \mu_{k,\ell}r_2) \cdot (I + \mu_{i,j} r_1)\cdot (I + \mu_{k,\ell}r_2).\\
  & = I + \mu_{i,j}(r_1 - r_1) + \mu_{k,\ell}(r_2 - r_2) = I\qedhere
  \end{align*}
\end{proof}

\noindent 
Therefore, for distinct $i,j,k \in [d]$ (which exist when $d\geq 3$), we have
\[
  \insquare{e_{i,j}(r_1), \insquare{e_{j,i}(r_2), e_{i,k}(r_3)}} = e_{i,k}(r_1r_2r_3)
\]
Thus, for all $d \geq 3$, using the above observation along with \expref{Observation}{obs:comm-generators}\eqref{item:comm-gen-sum}, we get that $e_{i,j}(r)$ for any $r\in R$ can be generated by $\setdef{e_{k,\ell}(at + b)}{k,\ell \in [d]\;,\; a,b\in \F_p\;}$. 
This in particular implies that $|G|$ is at least $p^s$.
On the other hand, the size of $K_i$ depends only on $d,p$ and is independent of $s$.
The lemma below describes $K_d$; the other %
$K_i$  %
are just rearrangements of rows and columns in $K_d$.

\begin{lemma}[Explicit description of $K_d$]
  The group $K_d = \inangle{e_{i,i+1}(at + b)\;:\; a,b\in \F_p\,,\, i\neq d}$ 
  consists of
 the %
%
matrices $A = (A_{i,j})$ of the following form:
  \begin{align*}
    A_{i,j} = \begin{cases}
      1 & \text{if }i = j,\\
      %
      %
%
      \text{a polynomial of degree $\leq j - i$} & \text{if }i < j,\\  %
      0 & \text{if }i > j.
    \end{cases}
  \end{align*}
  \noindent 
%
%
%
  Stating the above differently, %
  for any $n \in [d]$, the group 
 $K_n = \inangle{e_{i,i+1}(at + b)\;:\; a,b\in \F_p\,,\, i\neq n}$
%
%
%
consists of
the   %
matrices $A = (A_{i,j})$ of the following form:
  \begin{align*}
    A_{i,j} = \begin{cases}
      1 & \text{if }i = j,\\
      \text{a polynomial of degree $(j - i)\bmod{d}$} & \text{if }j\neq i \text{ and }n\not\in\set{i,i+1,\ldots,j-1}_{\bmod{d}},\\
      0 & \text{otherwise}.
    \end{cases}
  \end{align*}
\end{lemma}
%
%
\begin{proof}
%
Follows by repeated application  %
of \expref{Observation}{obs:comm-generators}. 
\end{proof}

\noindent
Therefore, we can obtain a crude bound of $|K_i| \leq p^{O(d^3)}$ for any $i$. Also, the above lemma also gives an explicit description of the groups $K_S$. 

\begin{corollary}[Explicit description of $K_S$]\label{cor:exp-desc-K_S}
For any subset $S \subseteq [d]$, the group $K_S = \Intersection_{i\in S} K_i$ consists of matrices $A = (A_{i,j})$ of the following form:
\begin{align*}
  A_{i,j} = \begin{cases}
    1 & \text{if }i = j,\\
    \text{a polynomial of degree $(j - i)\bmod{d}$} & \text{if }j\neq i \text{ and }\set{i,i+1,\ldots,j-1}_{\bmod{d}} \cap S = \emptyset,\\
    0 & \text{otherwise}.
  \end{cases}
\end{align*}
\end{corollary}

\noindent
Recall the \emph{other} set of subgroups defined for each $S \subseteq [d]$:
\[
  \widetilde{K_S} := \inangle{e_{i,i+1}(at+b)\;:\; a,b\in \F_p\,,\, i\notin S}.
\]
These groups can also be explicitly described.

\begin{lemma}[Explicit description of $\widetilde{K_S}$]\label{lem:exp-desc-K_S-tilde} For any $\emptyset \neq S \subseteq [d]$, the group  $\widetilde{K_S}$ is the set of all $d\times d$ matrices $A = (a_{ij})$ of the form
  \[
    a_{i,j} = \begin{cases}
      1 & \text{if }i = j,\\
      \text{a polynomial of degree $\leq j-i$} & \text{if }j\neq i \text{ and }\set{i,i+1,\ldots,j-1}_{\bmod{d}} \cap S = \emptyset,\\
      0 & \text{otherwise}.
    \end{cases}
  \]
\end{lemma}
\begin{proof}
Any $A\in \widetilde{K_S}$ can be expressed as  $A = B_1\cdots B_m$ where each $B_r = e_{i_{r},i_{r}+1}(\ell_r)$, for some linear polynomial $\ell_r$, with $i_r \notin S$.
Then,
\begin{align*}
  A_{i,j} & = \sum_{\substack{i_1,\ldots, i_{m+1}\\i_1 = i\;,\;i_{m+1} = j}} (B_1)_{i_1,i_2} (B_2)_{i_2,i_3} \cdots (B_m)_{i_m,i_{m+1}}. 
\end{align*}
From the structure of each $B_r$, any nonzero contribution from the RHS must involve either $i_{r+1} = i_{r}$, or $i_{r+1} = i_r + 1$ if $r \notin S$. This forces that the only entries of $A$ that are nonzero, besides the diagonal, are at $(i,j)$ with none of $\set{i,i+1,\ldots, j-1}$ in $S$.

In the case when $\set{i,i+1,\ldots, j-1} \cap S = \emptyset$, the above
argument also shows that the entry $A_{i,j}$ has degree at most $j-i$. Furthermore, using
\expref{Observation}{obs:comm-generators}, we can easily see that $e_{i,j}(f) \in \widetilde{K_S}$ for an arbitrary polynomial $f(t)$ of degree at most $j-i$. From this, we can deduce that the structure of $\widetilde{K_S}$ is exactly as claimed. 
\end{proof}

\begin{proof}[Proof of Lemma~\ref{lem:intersection-property-KS}.]
Follows immediately from \expref{Corollary}{cor:exp-desc-K_S} and \expref{Lemma}{lem:exp-desc-K_S-tilde}. 
\end{proof}

\noindent
From this point on, since the groups $K_S$ and $\widetilde{K_S}$ are identical, we drop the tilde notation and use $K_S$ for $\widetilde{K_S}$. 

\subsection{Connectivity of the coset complex}

\begin{lemma}\label{lem:KS-intersection-property}
  Let $S \subset [d]$ with $|S| \leq d-2$. Then,
  \[
    K_S= \inangle{K_{S} \cap K_i\colon i \in [d] \setminus
      S}. 
  \]
\end{lemma}
\begin{proof}
  It is clear that $K_S$ is a superset of the RHS. It only remains to show that the other containment also holds. To see this, consider an arbitrary generator $e_{j,j+1}(r)$ of $K_S$. Since $j \notin S$ and $|S| \leq d-2$, there is some $i \in [d] \setminus \inparen{S \union \set{j}}$. Therefore, $e_{j,j+1}(r) \in K_S \cap K_i$ and hence is generated by the RHS. 
\end{proof}

\noindent
Combining the above lemma with \expref{Lemma}{lem:conn-criterion} and \expref{Lemma}{lem:link-structure}, we have the following corollary.

\begin{corollary}
  For the coset complex $\mathcal{X}(G,\set{K_1,\ldots, K_d})$ defined by the above groups, the $1$-skeleton of every link is connected. 
\end{corollary}

\section{Spectral expansion of the complex}

In this section we prove that the coset complex $\X(G,\set{K_1,\ldots, K_d})$ 
is a good spectral HDX. The
%
Descent Theorem (\expref{Theorem}{thm:trickle-down})  %
states that it suffices to show that the $1$-dimensional links of 
faces in $\X(d-3)$ are good spectral expanders.

\subsection{Structure of $1$-dimensional links}

One-dimensional %
links of the coset complex constructed are links of $v
\in \mathcal{X}(G,\set{K_1,\ldots, K_d})$ of size exactly $d-2$ (which are elements of $\X(d-3)$). Any
such $v$ can be written as $\set{gK_1,\ldots, gK_d} \setminus
\set{gK_{i},gK_{j}}$ for $i,j\in [d]$ with $i \neq j$ and $g \in
G$. Since the link of $v$ is isomorphic to the link of $\set{K_1,\ldots, K_d} \setminus
\set{K_{i},K_{j}}$, we might as well assume that $g=\mathrm{id}$.
These happen to be of two types depending on whether $i$ and $j$ are
consecutive or not. 

\begin{observation}
  Consider $v = \set{K_1,\ldots, K_d} \setminus \set{K_i,K_j}$ where $i$ and $j$ are \emph{not} consecutive (\ie, $(i-j) \neq \pm 1 \bmod d$). Then the $1$-dimensional link of $v$ is a complete bipartite graph.
\end{observation}
\begin{proof}
  Note that since $j \neq i\pm 1$, we have $[e_{i,i+1}(r_1), e_{j,j+1}(r_2)] = \mathrm{id}$ by \expref{Observation}{obs:comm-generators}. Hence, these two elements commute.

  The link of $v$ corresponds to the coset complex $\mathcal{X}(H,\set{H_1,H_2})$ where
  \begin{align*}
    H & = K_{[d]\setminus\set{i,j}} =  \inangle{e_{i,i+1}(at+b),e_{j,j+1}(at+b) \;:\; a,b\in \F_p},\\
    H_1 & = K_{[d]\setminus\set{i}} =  \inangle{e_{i,i+1}(at+b) \;:\; a,b\in \F_p},\\
    H_2 & = K_{[d]\setminus\set{j}} =  \inangle{e_{j,j+1}(at+b) \;:\; a,b\in \F_p}.
  \end{align*}
  Thus, the groups $H_1$ and $H_2$ commute with each other and hence any element of $h\in H$ can be written as $h = g_1 \cdot g_2$ where $g_1 \in H_1$ and $g_2 \in H_2$. \expref{Observation}{obs:nonempty-coset-intersection} implies that the resulting graph is the complete bipartite graph.    
\end{proof}

The interesting case is when $v = \set{K_1,\ldots, K_d}\setminus\set{K_i,K_{i+1}}$. 
Without loss of generality, we may focus on the link of $v = \set{K_3,K_4,\ldots, K_d}$. This corresponds to the coset complex $\mathcal{X}(H,\set{H_1, H_2})$
where
\begin{align*}
  H & = K_{3,4,\ldots, d} = \inangle{e_{1,2}(at + b), e_{2,3}(at + b) \;:\; a,b \in \F_p},\\
  H_1 & = K_{2,3,4,\ldots, d} = \inangle{e_{1,2}(at + b) \;:\; a,b \in \F_p},\\
  H_2 & = K_{1,3,4,\ldots, d} = \inangle{e_{2,3}(at + b) \;:\; a,b \in \F_p}.
\end{align*}
Hence, it suffices to focus on the first three rows and columns of these matrices as the rest of them are constant. Written down explicitly, 
\begin{align*}
H & = \setdef{\begin{bmatrix}
  1 & \ell_1 & Q\\
  0 & 1 & \ell_2\\
  0 & 0 & 1
\end{bmatrix}}{\begin{matrix}\ell_1,\ell_2\text{ are linear polynomials in $\F_p[t]$ }\\\text{and $Q$ is a quadratic polynomial in $\F_p[t]$}
\end{matrix}},\\
  H_1 & = \setdef{\begin{bmatrix}
  1 & \ell & 0\\
  0 & 1 & 0\\
  0 & 0 & 1
\end{bmatrix}}{\ell\text{ is a linear polynomial in $\F_p[t]$}},\\
    H_2 & = \setdef{\begin{bmatrix}
  1 & 0 & 0\\
  0 & 1 & \ell\\
  0 & 0 & 1
\end{bmatrix}}{\ell\text{ is a linear polynomial in $\F_p[t]$}}.
\end{align*}
Multiplication of an arbitrary element of $H$ with an arbitrary element of $H_1$ is of the form
\[
\begin{bmatrix}
1 & \ell_1 & Q\\
0 & 1 & \ell_2\\
0 & 0 & 1
\end{bmatrix} \cdot \begin{bmatrix}
  1 & \ell & 0\\
  0 & 1    & 0\\
  0 & 0    & 1
  \end{bmatrix} = \begin{bmatrix}
    1 & \ell_1 + \ell & Q\\
    0 & 1 & \ell_2\\
    0 & 0 & 1
    \end{bmatrix}.
\]
Note that the a unique choice of $\ell$ that makes the $(1,2)$-th entry of the RHS zero is $\ell = -\ell_1$. Thus, each coset of $H_1$ in $H$ has a unique representative of the form $M_1(\ell,Q)$ described below, and similarly, each coset of $H_2$ %
in  %
%
$H$ has a unique representative of the form $M_2(\ell,Q)$: %
\[
M_1(\ell,Q) :=   \begin{bmatrix}
    1 & 0 & Q\\
    0 & 1 & \ell\\
    0 & 0 & 1
  \end{bmatrix} \quad,\quad M_2(\ell,Q) :=   \begin{bmatrix}
    1 & \ell & Q\\
    0 & 1 & 0\\
    0 & 0 & 1
  \end{bmatrix}
\]
respectively, %
where $\ell$ is a linear polynomial and $Q$ is a quadratic polynomial in $\F_p[t]$. %
%
This is because %
any element  %
of $H$ can be uniquely written as %
\begin{align*}
  \begin{bmatrix}
    1 & \ell_1 & Q\\
    0 & 1 & \ell_2\\
    0 & 0 & 1
  \end{bmatrix} & = \begin{bmatrix}
    1 & \ell_1 & Q - \ell_1\ell_2\\
    0 & 1 & 0\\
    0 & 0 & 1
  \end{bmatrix}\begin{bmatrix}
    1 & 0 & 0\\
    0 & 1 & \ell_2\\
    0 & 0 & 1
  \end{bmatrix}\\
  &= \begin{bmatrix}
    1 & 0 & Q - \ell_1\ell_2\\
    0 & 1 & \ell_2\\
    0 & 0 & 1
  \end{bmatrix}
  \begin{bmatrix}
    1 & \ell_1 & 0\\
    0 & 1 & 0\\
    0 & 0 & 1
  \end{bmatrix}.
\end{align*}

\begin{lemma}
  For linear polynomials $\ell_1,\ell_2 \in \F_p[t]$ and quadratic polynomials $Q_1,Q_2 \in \F_p[t]$, we have that
  \[
    M_1(\ell_1,Q_1) H_1 \cap M_2(\ell_2,Q_2) H_2 \neq \emptyset \quad\Longleftrightarrow\quad \ell_1\ell_2 = Q_1 - Q_2.
  \]
\end{lemma}
\begin{proof}
  Note that matrices in $H_1H_2$ are of the form
  \[
    \begin{bmatrix}
      1 & \ell_1 & 0\\
      0  & 1 & 0\\
      0  & 0 & 1
    \end{bmatrix}
    \begin{bmatrix}
      1 & 0 & 0\\
      0  & 1 & \ell_2\\
      0  & 0 & 1
    \end{bmatrix} = \begin{bmatrix}
      1 & \ell_1 & \ell_1\ell_2\\
      0  & 1 & \ell_2\\
      0  & 0 & 1
    \end{bmatrix}.
  \]
  By \expref{Observation}{obs:nonempty-coset-intersection}, the cosets have a non-empty intersection if and only if
  \[
    H_1 H_2 \ni M_1(\ell_1,Q_1)^{-1} M_2(\ell_2,Q_2) = \begin{bmatrix}
      1 & 0 & -Q_1\\
      0 & 1 & -\ell_1\\
      0 & 0 & 1
    \end{bmatrix}\begin{bmatrix}
      1 & \ell_2 & Q_2\\
      0 & 1 & 0\\
      0 & 0 & 1
    \end{bmatrix} = \begin{bmatrix}
      1 & \ell_2 & Q_2 - Q_1\\0 & 1 & -\ell_1 \\ 0 & 0 & 1
    \end{bmatrix}
  \]
  which happens if and only if $(\ell_2)(-\ell_1) = Q_2 - Q_1$ which is the same as $Q_1 - Q_2 = \ell_1\ell_2$. 
\end{proof}

Therefore, the $1$-dimensional link is the bipartite graph $A=(U,V,E)$
with left and right vertices identified by pairs $(\ell,Q)$ where
$\ell$ and $Q$ are linear and quadratic polynomials 
in $\F_p[t]$,  %
respectively, with
$(\ell_1,Q_1) \sim (\ell_2,Q_2) \Leftrightarrow \ell_1\ell_2 = Q_1 +
Q_2$ (by associating $M_1(\ell,Q)$ with the tuple $(\ell,Q)$ on the
left, and $M_2(\ell,Q)$ with the tuple $(\ell,-Q)$ on the right).

Note that $A$ is %
a  %
$p^2$-regular bipartite graph with  $p^5$ vertices on each side.
It suffices to show that $A$ is a good expander. 


Kaufman and Oppenheim~\cite{KaufmanO2018} prove the expansion
properties of this graph using
the  %
representation theory of the associated
groups, while we directly analyse the spectral gap of the adjacency
matrix associated with this graph. O'Donnell and Pratt~\cite[Case 2 in
the Proof of Theorem 3.23]{ODonnellP2022} give yet another proof of the 
spectral gap using the Polynomial Identity Lemma (also referred to as 
the Schwartz--Zippel lemma). 


\subsection{A related graph}

The following graph is the ``lines-points'' or the ``affine plane'' graph used by Reingold, Vadhan and Wigderson~\cite{ReingoldVW2005} (as the \emph{base graph} in
their  %
 construction of constant-degree expanders, using the zig-zag product). 
Let $\F_q$ be a finite field.
Consider the bipartite graph $B_q = (U',V', E')$ defined as follows:
\begin{align*}
%
%
%
U'=V' & = \F_q \times \F_q &E' & = \setdef{\inparen{(a,b),(c,d)}}{ac = b+d}.%
\end{align*}
Note that the graph $B_q$ is $q$-regular as for any 
%
%
$a,b,c \in \F_q\,$,  %
there is a unique $d \in \F_q$ such that $ac = b+d$ and thus the vertex $(a,b)$ has exactly $q$ neighbours in $B_q$. 
\begin{lemma}[{\cite[Lemma 5.1]{ReingoldVW2005}}]
  The $q$-regular bipartite graph $B_q$ is a
  $\frac{1}{\sqrt{q}}$-one-sided spectral expander.
\end{lemma}
\begin{proof}
  Let $B_q^2$ denote the multigraph whose adjacency matrix is the 
  square of the adjacency matrix of $B_q$. 
  %
  Note that $B_q^2$ is a multigraph as each edge in $B_q^2$ corresponds to
%
a length-two path  %
in $B_q$ and there may be more than one such path between a pair of vertices. 
%
%
%
It is easy to see that
\[
%
\text{Number of edges between $(a,b)$ and $(c,d)$ in $B_q^2$} =
\begin{cases} 1 & \text{if $a \neq c$,}\\
      q & \text{if $a = c$ and $b = d$,}\\
      0 & \text{otherwise.}
    \end{cases}
    \]
Therefore, the adjacency matrix of $B_q^2$
%
%
%
can be written\footnote{In the equation, the notation $\otimes$ refers to the
Kronecker product, or tensor product of matrices. It is well known
that, for square matrices $A$ and $B$, the
%
multiset  %
of eigenvalues of
%
$A \otimes B$  %
is all products of the form $\lambda_i \cdot \nu_j$ where $\lambda_i$ is an eigenvalue 
of $A$ and $\nu_j$ is an eigenvalue of $B$.} (under a suitable order of listing vertices) as
  \[
    q I_{q^2} + (J_q - I_q)\otimes J_q \quad\text{(where $J_q$ is the $q\times q$ matrix of $1s$)}.
  \]
  By observing that $J_q$ has eigenvalue of $q$ with multiplicity $1$, and eigenvalue $0$ with multiplicity $(q-1)$, a simple calculation shows that $B_q^2$ has eigenvalue of $q^2$ with multiplicity $1$, eigenvalue $q$ with multiplicity $q(q-1)$ and eigenvalue $0$ with multiplicity $q-1$. 
  Hence the unnormalized second largest eigenvalue of $B_q^2$ is $q$ and hence we have that the normalized second largest eigenvalue of $B_q$ is $1/\sqrt{q}$. 
\end{proof}

\subsection{Relating the graph $B_q$ with $A$}

Set $q = p^3$ so that %
$\F_q = \F_p[y]/\inangle{\mu(y)}$ %
for some irreducible polynomial $\mu(y)$ of %
%
degree~3.  %
Therefore, each
element in $\F_q$ is expressible as $a_0 + a_1y + a_2 y^2$ for some
$a_0,a_1,a_2\in \F_p$. Thus, the graph $B_q=(U',V',E')$ defined above,
for %
$q = p^3$, is a $p^3$-regular bipartite graph with $p^6$ vertices on either side.

Let $U'' = V'' = \setdef{(a_0 + a_1y , b_0 + b_1y + b_2 y^2)}{a_0,a_1,b_0,b_1,b_2 \in \F_p}$, which is a subset of $U'$ and $V'$, respectively, of size $p^5$ each.

\begin{observation}
  The induced subgraph of $B_q$ on  $U'',V''$ is exactly the graph
  $A=(U,V,E)$ described earlier. 
\end{observation}
\begin{proof}
  Note that $((\ell_1(y),Q_1(y)),(\ell_2(y),Q_2(y))) \in E'$ if and only if
  \[
    \ell_1(y)\cdot \ell_2(y) = Q_1(y) + Q_2(y) \bmod{\mu(y)}. 
  \]
  However, since the above equation has degree at most $2$, we have
  \[
    \ell_1(y)\cdot \ell_2(y) = Q_1(y) + Q_2(y) \quad\Leftrightarrow\quad \ell_1(y)\cdot \ell_2(y) = Q_1(y) + Q_2(y) \pmod{\mu(y)},
  \]
  and the first equation is exactly the adjacency condition of the graph $A$. Hence, the induced subgraph of $B_q$ on $U'',V''$ is indeed the graph $A$. 
\end{proof}

Normally, induced subgraphs of expanders need not even be connected. However, the following lemma shows that there are some instances where we may be able to give non-trivial bounds on $\lambda$.

\begin{lemma}\label{lem:induced-subgraph-eig-bound}
%
Let $X$ be  %
a $d$-regular %
graph that is an induced subgraph of a $D$-regular graph $Y$. Then,
  \[
    \lambda(X) \leq \frac{D \lambda(Y)}{d}. 
  \]
\end{lemma}
\begin{proof} The %
  Courant--Fischer  %
  characterization of the second largest eigenvalue
  (\expref{Theorem}{thm:cf1})  %
tells us that \(\lambda(X) = \max_{a \perp \mathbbm{1}_{|X|}} \frac{a^T G a}{d \cdot a^T a}\) where $G$ is the adjacency matrix of $X$.
 Consider an arbitrary $a \in \R^{|X|}$ such that $a \perp \mathbbm{1}_{|X|} = 0$. Since $X$ is an induced subgraph of $Y$, the vector $a$ can be padded with zeroes to obtain a vector $b_a \in \R^{|Y|}$ such that $b_a \perp \mathbbm{1}_{\abs{Y}}$. Therefore, if $A_X$ and $A_Y$ are the normalised adjacency matrices of $X$ and $Y$, we have
\[
   \lambda(X) = \max_{a \perp \mathbbm{1}_{|X|}} \frac{a^T A_X
   a}{a^T a} =  \frac{D}{d}\cdot \max_{a \perp \mathbbm{1}_{|X|}}
   \frac{b_a^T A_Y b_a}{b_a^T b_a}
   \leq \frac{D}{d}\cdot \max_{b \perp \mathbbm{1}_{|Y|}}
   \frac{b^T A_Y b}{b^T b} =
   \frac{D\lambda(Y)}{d} \enspace . \qedhere
 \]
\end{proof}

\begin{corollary}
  The graph $A(U,V,E)$ corresponding to the $1$-dimensional links of $\mathcal{X}(G,\set{K_1,\ldots, K_d})$ is a $\frac{1}{\sqrt{p}}$-one-sided spectral expander. 
\end{corollary}
\begin{proof}
  The graph $B_{p^3}$ is a bipartite, $p^3$-regular graph with $\lambda(B_{p^3}) \leq \frac{1}{p^{3/2}}$ and $A(U,V,E)$ is a $p^2$-regular graph that is an induced subgraph of $B_{p^3}$. Hence, by \expref{Lemma}{lem:induced-subgraph-eig-bound},
  \[
\lambda(A) \leq \frac{p^3 \cdot (1/p^{3/2})}{p^2} = \frac{1}{\sqrt{p}}. \qedhere
  \]
\end{proof}

\subsection*{The final expansion bounds}

From the corollary above, we obtain the following theorem of Kaufman and 
Oppenheim.

\begin{theorem}[\cite{KaufmanO2018}]\label{thm:KO-HDX} For $p > (d-2)^2$, 
  the $(d-1)$-dimensional coset complex $\X(G,\set{K_1,\ldots,K_d})$
  is a $\frac1{\sqrt{p}-(d-2)}$-one-sided spectral HDX.
\end{theorem}

\begin{proof}
%
It follows  %
directly from \expref{Theorem}{thm:trickle-down} that
  $\X(G,\set{K_1,\ldots,K_d})$ is a $\gamma$-one-sided spectral HDX for
  \[
  \gamma \leq \frac{\sfrac{1}{\sqrt{p}}}{1 - (d-2)(\sfrac{1}{\sqrt{p}})} = \frac{1}{\sqrt{p} - (d-2)}.\qedhere
  \]
\end{proof}

\paragraph{Constructing two-sided spectral HDXs and standard expanders:}

The $(d-1)$-dimensional coset complex $\X(G,\set{K_1,\ldots,K_d})$ is
not a two-sided spectral HDX as the 1-skeletons of the links of the faces in
$\X(d-3)$ are bipartite. However, if we restrict attention to the $k$-skeleton of
$\X$ for some $k<d-1$ then we can bound the least eigenvalue
%
by applying the Descent Theorem to the least eigenvalue %
(\expref{Theorem}{thm:trickle-down-one-level}\eqref{thm:trickle-down-negative}).
This is summarized in the following corollary. 

\begin{corollary}\label{cor:KO-exp} For $p > (d-2)^2$ and any $1 \leq
  k < d$
  the $k$-skeleton of the $(d-1)$-dimensional coset complex $\X(G,\set{K_1,\ldots,K_d})$
is a $\max\set{\frac1{\sqrt{p}-(d-2)}\;,\;\frac{1}{d-k}}$-two-sided spectral HDX.
  \end{corollary}

  In particular, if we set $k=1$ in the above corollary, we get a
  standard
  $\max\set{\frac1{\sqrt{p}-(d-2)}\;,\;\frac{1}{d-1}}$-two-sided spectral expander.  This graph is a $d$-partite graph and hence its
  least eigenvalue is at most $\sfrac{-1}{(d-1)}$, while the above
  argument shows that it is least (and hence equal to)
  $\sfrac{-1}{(d-1)}$. 

 Thus, this not only yields an elementary construction and proof of
  one-sided spectral HDXs (\expref{Theorem}{thm:KO-HDX}), but also one of
  standard spectral expander (\expref{Corollary}{cor:KO-exp}).
  
  


\appendix


\section{Proof of the Descent Theorem}\label{sec:descent}

For the sake of completeness, we present the proof of
\expref{Theorem}{thm:trickle-down} that asserts that proving spectral expansion
for the maximal faces is sufficient to obtain expansion of any
link. This exposition is essentially from the lecture notes by 
Dikstein~\cite{Dikstein2019}. 

Let $(X,w)$ be a weighted $d$-dimensional simplicial complex. Let
$\mu_{d}=w|_{X(d)}$ be the distribution on the set $X(d)$ of
$(d+1)$-sized faces. This distribution induces distributions $\mu_i$ on
$X(i)$ in the natural way.\\ 

For two functions $f,g\colon X(0) \rightarrow \R$, define their \emph{inner product}
$\inangle{f,g}_X = \E_{u \sim \mu_0} [f(u) g(u)]$. We will drop the
subscript $X$ if it is clear from context. Note that, by the
definition of $\mu_1$, sampling $u$ according to $\mu_0$  can be
equivalently achieved  by sampling an edge $(u,v)$ according to $\mu_1$ and returning one of the points uniformly at random. Therefore,
\begin{equation}\label{eqn-innerprod1}
  \inangle{f,g}_X = \E_{u \sim \mu_0}[f(u) g(u)] = \E_{\set{u,v} \sim
    \mu_1}[f(u)g(u)] =\E_{v \sim \mu_0} \E_{u \sim X_v(0)}[f(u) g(u)] =\E_{v \sim \mu_0}[\inangle{f_v,g_v}_{X_v}],
\end{equation}
where $f_v,g_v\colon X_v(0) \rightarrow \R$ are the restrictions to the link of $v$. 

Define the \emph{adjacency operator} $A$ that, on a function $f\colon X(0) \rightarrow \R$ on vertices returns another function $Af$ on vertices defined via
\[
  Af(v) = \E_{u\sim v}[f(u)],
\]
where $u\sim v$ refers to a random neighbour of $v$ according to the distribution $u \sim \mu_0(X_v)$.
In other words, $A$ \emph{averages} $f$ over neighbours. Furthermore,
$A$ is self-adjoint with respect to the above inner product, i.e,
$\inangle{Af,g} = \inangle{f,Ag}$. Hence, it has $n$ real eigenvalues
and an orthonormal set of eigenvectors. Clearly
$A\mathbbm{1} = \mathbbm{1}$; the constant $1$ function is an
eigenvector for this operator (in fact, it is an eigenvector
corresponding to the largest eigenvalue 1). The remaining eigenvalues
are characterized by the %
Courant--Fischer Theorem, stated below.  %

%
\begin{theorem}[Courant--Fischer]  \label{thm:cf1}  %
  Let $A\in \R^{n \times n}$ be an $n \times n$ matrix over the reals
  that is self-adjoint with respect to some inner product
  $\inangle{\cdot , \cdot }\colon \R^n \times \R^n \to \R$. Then $A$ has $n$ real
  eigenvalues $\lambda_1 \geq \cdots \geq \lambda_n$ which have the
  following characterization.
  \[
    \lambda_i = \max_{V \colon \dim V = i}\; \min_{0\neq  x \in V}
    \frac{\inangle{x,Ax}}{\inangle{x,x}} =  \min_{V \colon \dim V = n-i+1}\; \max_{ 0\neq x \in V}
    \frac{\inangle{x,Ax}}{\inangle{x,x}}.\qedhere
  \]
%
\end{theorem}       %
    


Similar to \eqref{eqn-innerprod1}, we have
\begin{align}
  \inangle{Af,g}_X & = \E_{\set{u,w} \sim \mu_1} [f(u) g(w)] = \E_{\set{u,v,w} \sim \mu_2} [f(u) g(w)]\nonumber\\
                 & = \E_{v \sim \mu_0}\insquare{\E_{\set{u,w}\sim  \mu_1(X_v)} [f(u) g(w)]}\nonumber\\
                 & = \E_{v \sim \mu_0}\insquare{\inangle{A_v f_v,g_v}_{X_v}}\label{eqn-innerprod2}
\end{align}
where $A_v$ denotes the adjacency operator restricted to the link $X_v$. \\

With the above notation, we can now state the theorem we wish to prove.
It suffices to prove the theorem in the case of $d = 3$ as we can obtain 
\expref{Theorem}{thm:trickle-down} by induction.

\begin{theorem}\label{thm:trickle-down-one-level}
  Suppose $(X,w)$ is weighted $2$-dimensional simplicial complex. Then, we have the following two implications:
  \begin{enumerate}
  \item \label{thm:trickle-down-positive} Suppose the 1-skeleton of
    $X$ is connected and for every vertex $v \in X(0)$,
    $\inangle{A_v f,f} \leq \lambda \inangle{f,f}$ for all $f\colon
    X_v(0) \rightarrow \R$ with $f \perp \mathbbm{1}_{X_v}$ for some
    $\lambda \in [0,1)$. Then, for
    any $g\colon X(0) \rightarrow \R$ with $g \perp \mathbbm{1}_X$, we
have $\inangle{Ag,g} \leq \gamma \inangle{g,g}$ where $\gamma \leq
\frac{\lambda}{1- \lambda}$.
  \item \label{thm:trickle-down-negative} Suppose the 1-skeleton of
    $X$ is non-empty and for every vertex $v \in X(0)$, we have
    $\inangle{A_v f,f} \geq \eta \inangle{f,f}$ for all $f\colon
    X_v(0) \rightarrow \R$ for some $\eta \in [-1,1)$. Then, for any $g\colon X(0) \rightarrow \R$, we
have $\inangle{Ag,g} \geq \gamma \inangle{g,g}$ where $\gamma \geq
\frac{\eta}{1- \eta}$.
\end{enumerate}
  

\end{theorem}

\noindent Before we see a proof of this, let us  see how \expref{Theorem}{thm:trickle-down} follows from this.


\begin{dtheorem}[\expref{Theorem}{thm:trickle-down} restated]\label{thm:td-full}
  Suppose $(X,w)$ is a non-empty $d$-dimensional weighted simplicial complex with the following properties.
  \begin{itemize}\itemsep0pt
  \item The $1$-skeleton of every link is connected.
  \item For all $v \in X(d-2)$, the link $(X_v,w_v)$ is a
    $\lambda$-one-sided spectral expander for some $\lambda < \frac1{d-1}$. \Ie, there is a $\lambda > 0$ such that, for every $v \in X(d-2)$ and every $g\colon X_v(0) \rightarrow \R$ with $g \perp \mathbbm{1}$, we have
  \[
\inangle{A_v g,g} \leq \lambda \inangle{g,g}.
  \]

  \end{itemize}
Then, $(X,w)$ is a $\gamma$-one-sided spectral HDX for $\gamma \leq \frac{\lambda}{1 - (d-1)\lambda}$. That is, for any $v \in X(-1) \union \cdots \union X(d-2)$ and every $g\colon X_v(0) \rightarrow \R$ with $g \perp \mathbbm{1}$, we have $\inangle{A_v g,g} \leq \gamma \inangle{g,g}$.\\
  

Furthermore, suppose we also know that there is a $\eta \in [-1,0)$ such that, for every $v \in X(d-2)$ and every $g\colon X_v(0) \rightarrow \R$, we have $\inangle{A_v g,g} \geq \eta \inangle{g,g}$. Then, $X$ is a $\gamma$-two-sided spectral HDX with 
\[
  \gamma \leq \max\inparen{\frac{\lambda}{1 - (d-1)\lambda}, \abs{\frac{\eta}{1 - (d-1)\eta}}}.
\]
That is, for every $g \colon X_v(0) \rightarrow \mathbb{R}$ with $g
\perp \mathbbm{1}$, we have $|\inangle{A_v g,g} | \leq \gamma \inangle{g,g}$.
\end{dtheorem}

\begin{proof}
  For any $i \leq d-2$, let
  \[\lambda_i = \min_{v\in X(i)} \max_{\substack{g:X_v(0) \rightarrow \R\\g\perp \mathbbm{1}}}\frac{\inangle{A_v g,g}}{\inangle{g,g}},
  \]
  the smallest one-sided spectral expansion with respect to $X(i)$. From repeated applications of  \expref{Theorem}{thm:trickle-down-one-level},
  \[
    \lambda_{-1} \leq \frac{\lambda_0}{1 - \lambda_0} \leq \frac{\lambda_{1}/(1-\lambda_{1})}{1 - (\lambda_{1}/(1-\lambda_{1}))} = \frac{\lambda_{1}}{1 - 2\lambda_{1}} \leq \cdots \leq \frac{\lambda_{d-2}}{1 - (d-1)\lambda_{d-2}}
  \]
  which eventually completes the proof for one-sided spectral expansion.

  For two-sided spectral expansion, we also have to show that all the eigenvalues are bounded away from $-1$. One again, let $\eta_i$ be such that
  \[
    \eta_i = \max_{v\in X(i)} \min_{\substack{g\colon X_v(0) \rightarrow \R\\g\perp \mathbbm{1}}}\frac{\inangle{A_v g,g}}{\inangle{g,g}}.    
  \]
  By repeated applications of \expref{Theorem}{thm:trickle-down-one-level}~\eqref{thm:trickle-down-negative}, we obtain
  \[
    \eta_{-1} \geq \frac{\eta_0}{1 - \eta_0} \geq \frac{\eta_{1}/(1-\eta_{1})}{1 - (\eta_{1}/(1-\eta_{1}))} = \frac{\eta_{1}}{1 - 2\eta_{1}} \geq \cdots \geq \frac{\eta_{d-2}}{1 - (d-1)\eta_{d-2}}
  \]
  Together, we have that $X$ is a $\gamma$-two-sided spectral HDX for
  \[
  \gamma = \max\inparen{\frac{\lambda}{1 - (d-1)\lambda}, \abs{\frac{\eta}{1 - (d-1)\eta}}}.\qedhere
  \]
  \end{proof}


\begin{proof}[Proof of Theorem~\ref{thm:trickle-down-one-level}]
  Let $g$ be an eigenvector that satisfies $\inangle{g,g} = 1$ and
  $g\perp \mathbbm{1}_X$ that maximises (or minimises) $\inangle{Ag,g}$, and $\gamma =
  \inangle{Ag,g}$ be the extremal value. In particular, $Ag =
  \gamma\cdot g$. From \eqref{eqn-innerprod2} we have
    $\gamma = \inangle{Ag,g} = \E_v \insquare{\inangle{A_v g_v,g_v}}$.

    Even though $g \perp \mathbbm{1}_X$, the \emph{local} component $g_v$ need not be perpendicular to $\mathbbm{1}_{X_v}$. Hence,  let us write $g_v = \alpha_v \mathbbm{1}_{X_v} + g_v^{\perp}$ where $g_v^{\perp} \perp \mathbbm{1}_{X_v}$; we shall drop the subscript from $\mathbbm{1}_{X_v}$ for the sake of brevity as the length of the vector will be clear from context. Note that $\alpha_v = \inangle{g_v, \mathbbm{1}} = \E_{w\in X_v(0)} [g_v] = Ag(v)$. Therefore, $\E_v [\alpha_v^2] = \inangle{Ag,Ag} = \gamma^2$. Hence,
  \begin{align}
    \gamma = \inangle{Ag,g} & = \E_v \insquare{\inangle{A_v g_v,g_v}} = \E_v \insquare{\alpha_v^2 + \inangle{A_v g_v^\perp, g_v^{\perp}}}\label{eqn:trickle-LHS}
  \end{align}
  We shall now focus on the proof of \expref{Theorem}{thm:trickle-down-one-level}~\eqref{thm:trickle-down-positive}. The other direction is exactly identical with the inequality flipped.

  In the case of \expref{Theorem}{thm:trickle-down-one-level}~\eqref{thm:trickle-down-positive}, where we are given $\inangle{A_v g_v^\perp, g_v^{\perp}} \leq \lambda \inangle{g_v^\perp, g_v^{\perp}}$ for all $v\in X(0)$, we have
  \begin{align*}
    \gamma &= \E_v \insquare{\alpha_v^2 + \inangle{A_v g_v^\perp, g_v^{\perp}}} \leq \E_v \insquare{\alpha_v^2 + \lambda \inangle{g_v^\perp, g_v^{\perp}}} \\
    & = \E_v \insquare{(1-\lambda) \alpha_v^2 + \lambda \inangle{g_v, g_v}}\\
    &= (1-\lambda) \gamma^2 + \lambda.\\
    \implies \gamma(1-\gamma) & \leq \lambda (1- \gamma^2)\\
    \implies \gamma & \leq \lambda (1+\gamma) &\text{(connected, thus $\gamma < 1$)}\\
    \implies \gamma & \leq \frac{\lambda}{1 - \lambda}.
  \end{align*}

  In the case of \expref{Theorem}{thm:trickle-down-one-level}~\eqref{thm:trickle-down-negative}, where we are given $\inangle{A_v g_v^\perp, g_v^{\perp}} \geq \eta \inangle{g_v^\perp, g_v^{\perp}}$ for all $v\in X(0)$, the same argument yields
  \begin{align*}
    \gamma = \E_v \insquare{\alpha_v^2 + \inangle{A_v g_v^\perp, g_v^{\perp}}}&\geq \E_v \insquare{\alpha_v^2 + \eta \inangle{g_v^\perp, g_v^{\perp}}} = (1 - \eta)\gamma^2 + \eta\\
    \implies \gamma & \geq \frac{\eta}{1 - \eta}\qedhere
  \end{align*}
\end{proof}



\section{Primer on group theory}\label{sec:group}

In this section, for completeness,  we shall note the basic definitions and properties of groups that are used in this exposition. 

\begin{definition}[Groups and subgroups]
  A set of elements $G$ equipped with a binary operation $\star: G \times G \rightarrow G$ is said to be a \emph{group} if it satisfies the following properties:
  \begin{description}
  \item[Associativity:] For all $g_1,g_2,g_3 \in G$, we have $\inparen{g_1 \star g_2} \star g_3 = g_1 \star \inparen{g_2 \star g_3}$.
  \item[Identity:] There exists an \emph{identity} element $\mathrm{id} \in G$ such that, for all $g\in G$, we have $g \star \mathrm{id} = \mathrm{id} \star g = g$.
  \item[Inverses:] For every element $g \in G$, there is an element $g^{-1} \in G$ such that $g \star g^{-1} = g^{-1} \star g = \mathrm{id}$. 
  \end{description}
  
  
  A subset $H \subseteq G$ is said to be a \emph{subgroup} of $G$ if
  $H$ the binary operation $\star$ restricted to $H$ satisfies  the
  above three properties (including the fact that $h_1 \star h_2 \in
  H$ for all $h_1,h_2 \in H$).

   \medskip

  Often the binary operation $\star$ is omitted and products just expressed as concatenation of elements.

 
\end{definition}

\begin{definition}[Cosets]
Given a subgroup $H$ of a group $G$, if $x \in G$ is an arbitrary element, the 
%
\emph{left coset}      %
of $H$ containing $x$, denoted by $x H$, is defined as the set
\[
x H = \setdef{x h}{h\in H}.
\]
Two cosets $x H$ and $y H$ are identical if and only if $x^{-1} y \in H$. Hence,
any element $x'\in x H$ is also referred to as a \emph{coset representative} of
$x H$ as $x' H = x H$. 
\end{definition}

%
Right cosets  %
are defined similarly. A subgroup $H$ is said to be
\emph{normal} if the
%
right cosets and the left cosets  %
agree for all $x$,
\ie, $xH = Hx, \forall x \in G$. 

Since two cosets of a subgroup $H$ of $G$ are either identical or
disjoint, the set of distinct cosets of a subgroup $H$ of $G$
partition the elements of $G$. If a subgroup $H$ is normal, this set
of cosets forms a group $G/H$, called the \emph{quotient group} of $H$
in $G$.\\


%
For subgroups $H, K$ %
of $G$, we will often consider the product $HK$ (or $H \star K$) which refers to the set $\setdef{h k}{h \in H\;,\; k\in K}$. 
It is worth stressing that $HK$ \emph{need not} be a subgroup of $G$ and the above just refers to a set of elements that can be expressed as an (ordered) product of an element in $H$ and an element in $K$. 

For an arbitrary %
subset  %
$S$ of $G$, we will define $\inangle{S}$ as the smallest subgroup of $G$ that contains the set $S$. This is also referred to as the \emph{group generated} by $S$.\\


In general, the binary operation $\star$ is order dependent. Groups where $g_1 g_2 = g_2 g_1$ for all $g_1,g_2 \in G$ are said to be \emph{commutative} or \emph{abelian} groups. The following notion of \emph{commutators} (and \emph{commutator subgroups}) is a way to measure \emph{how non-commutative} a group $G$ is.

\begin{definition}[Commutators]\label{defn:commutator}
  For a pair of elements $g, h\in G$, we shall define the \emph{commutator} of $g,h$ (denoted by $[g,h]$) as
  \[
    [g,h] := g^{-1} h^{-1} g h. 
  \]
  The \emph{commutator subgroup} of $G$, denoted by $[G,G]$ is the \emph{group generated by all commutators}. That is,
  \[
    [G,G] := \inangle{\setdef{[g,h]}{g,h\in G}}.\qedhere
  \]
\end{definition}

Note that if $G$ is abelian, then $[G,G] = \set{\mathrm{id}}$. As mentioned earlier, the commutator subgroup can be thought of as a way of describing how non-abelian a group is. In fact, the commutator subgroup of $G$ is the smallest \emph{normal} subgroup $H$ of $G$ such that the \emph{quotient} $G/H$ is abelian (although these are concepts that are not necessary to follow this exposition).



%% Theory of Computing
%% File: v020a005.bbl
%% Publication date: October 26, 2024
%% Authors: Prahladh Harsha and Ramprasad Saptharishi
\newcommand{\prelim}{Preliminary version in\ }\newcommand{\prelims}{Preliminary
  versions in\ }\newcommand{\sortsilent}[1]{}
% Define empty bibhead if not already defined
\providecommand{\bibhead}[1]{}
% Define tocrefpdfbookmark if not already defined
\expandafter\ifx\csname pdfbookmark\endcsname\relax%
  \providecommand{\tocrefpdfbookmark}{}
\else\providecommand{\tocrefpdfbookmark}{%
   \hypertarget{tocreferences}{}%
   \pdfbookmark[1]{References}{tocreferences}}%
\fi

\tocrefpdfbookmark
\begin{thebibliography}{10}

\bibitem{Dikstein2019}\bibhead{Dikstein2019}
{\sc Yotam Dikstein}: {Oppenheim's Trickling Down Theorem}, 2019.
\newblock {Lecture notes from the Error-Correcting Codes and High-Dimensional
  Expansion Boot Camp at Simons Institute for the Theory of Computing},
  Available at
  \href{https://simons.berkeley.edu/sites/default/files/docs/14119/ecctalkii.pdf}{https://simons.berkeley.edu/sites/default/files/docs/14119/ecctalkii.pdf}.

\bibitem{Garland1973}\bibhead{Garland1973}
{\sc Howard Garland}: $p$-adic curvature and the cohomology of discrete
  subgroups of $p$-adic groups.
\newblock {\em Ann. Math.}, 97(3):375--423, 1973.
\newblock [\epfmtdoi{10.2307/1970829}]

\bibitem{HM89}\bibhead{HM89}
{\sc Alexander~J. Hahn and O.~Timothy O'Meara}: {\em The {C}lassical {G}roups
  and {K}-Theory}.
\newblock Springer, 1989.
\newblock [\epfmtdoi{10.1007/978-3-662-13152-7}]

\bibitem{KaufmanO2018}\bibhead{KaufmanO2018}
{\sc Tali Kaufman and Izhar Oppenheim}: High dimensional expanders and coset
  geometries.
\newblock {\em Europ. J. Combinat.}, 111(103696), 2023.
\newblock \prelim \href{https://doi.org/10.1145/3188745.3188782}{STOC'18}.
\newblock [\epfmtdoi{10.1016/j.ejc.2023.103696}, \epfmt{arxiv}{1710.05304}]

\bibitem{LubotzkyPS1988}\bibhead{LubotzkyPS1988}
{\sc Alexander Lubotzky, Ralph Phillips, and Peter Sarnak}: Ramanujan graphs.
\newblock {\em Combinatorica}, 8(3):261--277, 1988.
\newblock \prelim \href{https://doi.org/10.1145/12130.12154}{STOC'86}.
\newblock [\epfmtdoi{10.1007/BF02126799}]

\bibitem{LubotzkySV2005-exphdx}\bibhead{LubotzkySV2005-exphdx}
{\sc Alexander Lubotzky, Beth Samuels, and Uzi Vishne}: Explicit constructions
  of {R}amanujan complexes of type $\tilde{A_d}$.
\newblock {\em Europ. J. Combinat.}, 26(6):965--993, 2005.
\newblock [\epfmtdoi{10.1016/j.ejc.2004.06.007}, \epfmt{arxiv}{math/0406217}]

\bibitem{LubotzkySV2005-hdx}\bibhead{LubotzkySV2005-hdx}
{\sc Alexander Lubotzky, Beth Samuels, and Uzi Vishne}: Ramanujan complexes of
  type $\tilde{A_d}$.
\newblock {\em Israel J. Math.}, 149(1):267--299, 2005.
\newblock [\epfmtdoi{10.1007/BF02772543}, \epfmt{arxiv}{math/0406208}]

\bibitem{margulis88}\bibhead{margulis88}
{\sc {\sortsilent{Margulis}{Grigorii} Aleksandrovich Margulis
  (\foreignlanguage{russian}{Григорий Александрович
  Маргулис})}}: \foreignlanguage{russian}{Явные
  теоретико-групповые конструкции
  комбинаторных схем и их применения в
  построении расширителей и
  концентраторов} ({R}ussian) [{E}xplicit group-theoretical
  constructions of combinatorial schemes and their application to the design of
  expanders and concentrators].
\newblock {\em Problemy Peredachi Informatsii}, 24(1):51--60, 1988.
\newblock Available at
  \href{https://www.mathnet.ru/eng/ppi686}{https://www.mathnet.ru/eng/ppi686},
  (English translation in \emph{Problems Inform. Transmission}, 24(1):39--46,
  1988).

\bibitem{ODonnellP2022}\bibhead{ODonnellP2022}
{\sc Ryan O'Donnell and Kevin Pratt}: High-dimensional expanders from
  {C}hevalley groups.
\newblock In {\em Proc. 37th Comput. Complexity Conf. (CCC'22)}, pp. 18:1--26.
  Schloss Dagstuhl--Leibniz-Zentrum fuer Informatik, 2022.
\newblock [\epfmtdoi{10.4230/LIPIcs.CCC.2022.18}, \epfmt{arxiv}{2203.03705}]

\bibitem{Oppenheim2018}\bibhead{Oppenheim2018}
{\sc Izhar Oppenheim}: Local spectral expansion approach to high dimensional
  expanders {P}art {I}: {D}escent of spectral gaps.
\newblock {\em Discr. Comput. Geom.}, 59(2):293--330, 2018.
\newblock [\epfmtdoi{10.1007/s00454-017-9948-x}, \epfmt{arxiv}{1709.04431}]

\bibitem{ReingoldVW2005}\bibhead{ReingoldVW2005}
{\sc Omer Reingold, Salil Vadhan, and Avi Wigderson}: Entropy waves, the
  zig-zag graph product, and new constant-degree expanders.
\newblock {\em Ann. Math.}, 155(1):157--187, 2005.
\newblock \prelim \href{https://doi.org/10.1109/SFCS.2000.892006}{FOCS'00}.
\newblock [\epfmtdoi{10.2307/3062153}, \epfmt{arxiv}{math.CO/0406038},
  \epfmt{eccc}{TR01-018}]

\bibitem{Sarveniazi2007}\bibhead{Sarveniazi2007}
{\sc Alireza Sarveniazi}: Explicit construction of a {R}amanujan
  $(n_1,n_2,\ldots,n_{d-1})$ -regular hypergraph.
\newblock {\em Duke Math. J.}, 136(1):141--171, 2007.
\newblock [\epfmtdoi{10.1215/S0012-7094-07-13913-9}]

\end{thebibliography}

\end{document}
